\section*{अध्याय \ref{ch:b3:mantle-convection} --- प्रावार संवहन और प्लेट-चालक बल}

\begin{solution}{exo:b3:mantle-convection:1}
$\mathrm{Ra} = 3500\times9.9\times\num{3e-5}\times1000\times(\num{6.7e5})^3/
(10^{-6}\times\num{5e20}) = \num{6.25e5}$, लगभग 950 गुना क्रांतिक
मान 657.5: परत प्रबल रूप से संवहन करती है।
\end{solution}

\begin{solution}{exo:b3:mantle-convection:2}
$\tau_M = 10^{21}/\num{7.7e10} = \qty{1.30e10}{s}$, लगभग \qty{412}{yr}।
ज्वार, जो घंटों तक चलता है, प्रत्यास्थ प्रावार का सामना करता है; हिमचादर का
हज़ारों वर्षों में होने वाला पिघलना, और लाखों वर्षों में होने वाला स्लैब का धँसना,
श्यान प्रावार का सामना करते हैं।
\end{solution}

\begin{solution}{exo:b3:mantle-convection:3}
$\tau = 4\pi\times10^{21}/(3300\times9.8\lambda)$: \qty{3.89e11}{s} =
\qty{1000}{km} के लिए \qty{12.3}{kyr}, और \qty{3000}{km} के लिए \qty{4.1}{kyr}।
व्यापक भार अधिक तेज़ी से उबरता है।
\end{solution}

\begin{solution}{exo:b3:mantle-convection:4}
$g\rho_m\alpha(T_m - T_0)\kappa t = 9.8\times3300\times\num{3.1e-5}\times1450
\times\num{0.8e-6}\times50\times\num{3.156e13} = \qty{1.835e12}{N.m^{-1}}$; कोष्ठक का
मान $1 + 2\times3300\times\num{3.1e-5}\times1450/(\pi\times2270) = 1.0416$,
अतः $F_{rp} = \qty{1.91e12}{N.m^{-1}}$, स्लैब-खिंचाव से 24 गुना कम है।
\end{solution}

\begin{solution}{exo:b3:mantle-convection:5}
(क) $v = 3300\times9.8\times\num{3e-5}\times300\times(\num{e4})^2/10^{21} =
\qty{2.91e-11}{m.s^{-1}}$; \qty{10}{km} के आर-पार ऊपर उठने का समय $d/v =
\qty{3.44e14}{s}$ और \qty{10.9}{Myr} है; विसरण-काल $d^2/\kappa =
\qty{e14}{s}$ और \qty{3.2}{Myr} है। इनका अनुपात $(d^2/\kappa)/(d/v) =
vd/\kappa = \rho g\alpha\Delta Td^3/(\kappa\eta) = \mathrm{Ra} = 0.29$ है। (ख) यह
पहले ठंडा होता है: छोटे पिंड संवहन नहीं करते। (ग) $\mathrm{Ra} \propto d^3$, अतः $d =
10\times(657.5/0.291)^{1/3} = \qty{131}{km}$ है।
\end{solution}

\begin{solution}{exo:b3:mantle-convection:6}
$2\mathrm{Ra}_c = 1315$ और $\mathrm{Ra}/2\mathrm{Ra}_c = 7604$, जिसका घनमूल
19.66 है। $\delta = 2891/19.66 = \qty{147}{km}$, $\mathrm{Nu} = 19.66/2 =
9.8$, $u = (10^{-6}/\num{2.891e6})\times19.66^2 = \qty{1.34e-10}{m.s^{-1}} =
\qty{0.42}{cm/yr}$, और $q = 9.8\times3.1\times1450/\num{2.891e6} =
\qty{15}{mW.m^{-2}}$ है।
\end{solution}

\begin{solution}{exo:b3:mantle-convection:7}
$u = 0.09/\num{3.156e7} = \qty{2.85e-9}{m.s^{-1}}$ और $B = 2270\times\num{8e8}
\times\num{2.85e-9} = \qty{5.18e3}{kg.s^{-1}} = \qty{5.2}{Mg/s}$, लगभग
हवाई के तीन-पाँचवें भाग के बराबर है। ऊष्मा-फ़्लक्स $1170\times5178/\num{3.1e-5} =
\qty{1.95e11}{W}$, \qty{0.2}{TW}, या पृथ्वी की ऊष्मा-हानि का
\qty{0.42}{\percent} भाग।
\end{solution}

\begin{solution}{exo:b3:mantle-convection:8}
$v = 2\times30\times9.8\times(\num{4e5})^2/(9\times10^{21}) =
\qty{1.05e-8}{m.s^{-1}} = \qty{0.33}{m/yr}$; \qty{2000}{km} तय करने में
$\num{2e6}/\num{1.05e-8} = \qty{1.9e14}{s}$ लगते हैं, लगभग \qty{6}{Myr}।
\end{solution}

\begin{solution}{exo:b3:mantle-convection:9}
$w = 300e^{-10/4.4} = \qty{30.9}{m}$ और $\dot w = w/\tau = 30.9/4400 =
\qty{7.0}{mm/yr}$, उमेओ के निकट मापी गई दर का दो-तिहाई:
इस $\tau$ वाला एक समान अर्ध-समष्टि सही परास में है, परंतु वास्तविक प्रतिक्षेप
अधिक तेज़ है या उसका अवनमन अधिक गहरा है।
\end{solution}

\begin{solution}{exo:b3:mantle-convection:10}
(क) जहाँ $w = 2\rho_m\alpha\Delta T\sqrt{\kappa t/\pi}/(\rho_m - \rho_w)$,
\[
  \tfrac12(\rho_m - \rho_w)gw^2 = \frac{2g\rho_m^2\alpha^2\Delta T^2\kappa t}
  {\pi(\rho_m - \rho_w)} = g\rho_m\alpha\Delta T\kappa t\times
  \frac{2\rho_m\alpha\Delta T}{\pi(\rho_m - \rho_w)} .
\] (ख) खंडशः समाकलन से,
$\int_0^\infty x\operatorname{erfc}x\,\mathrm dx = [\tfrac{x^2}2\operatorname{erfc}
x]_0^\infty + \int_0^\infty\tfrac{x^2}2\,\tfrac2{\sqrt\pi}e^{-x^2}\,\mathrm dx =
0 + \tfrac1{\sqrt\pi}\cdot\tfrac{\sqrt\pi}4 = \tfrac14$। (ग) \qty{100}{Myr} पर,
$\rho_mg\alpha\int z\theta = \qty{3.67e12}{N.m^{-1}}$ और $w =
\qty{3705}{m}$ लेने पर, $\tfrac12(\rho_m - \rho_w)gw^2 = \qty{1.53e11}{N.m^{-1}}$: जल-पद
केवल बल का \qty{4}{\percent} है।
\end{solution}

\begin{solution}{exo:b3:mantle-convection:11}
$\ln f = 3\ln(\pi^2 + k^2) - 2\ln k$, जिसका अवकलज $6k/(\pi^2 + k^2) - 2/k$
शून्य होता है जब $3k^2 = \pi^2 + k^2$, $k^2 = \pi^2/2$। वहाँ $f = (3\pi^2/2)^3/
(\pi^2/2) = 27\pi^4/4$। तरंगदैर्घ्य $2\pi/k = 2\sqrt2\,d = 2.83\,d$ प्रति-घूर्णी
संवहन-कोषों का एक युग्म धारण करती है, जिनमें से प्रत्येक $\sqrt2\,d = 1.41\,d$ चौड़ा है।
\end{solution}

\begin{solution}{exo:b3:mantle-convection:12}
(क) $u = \qty{1.90e-9}{m.s^{-1}}$, $\eta u/h = 10^{20}\times\num{1.90e-9}/
\num{1.5e5} = \qty{1.27}{MPa}$, और \qty{6000}{km} पर \qty{7.6e12}{N.m^{-1}}।
(ख) कटक-धक्का का दोगुना और स्लैब-खिंचाव का छठा भाग: \omterm{def:b3:mantle-convection:basal-drag}{आधारी कर्षण}
कटक-धक्के को अकेले रोक सकता है, पर स्लैब को नहीं। (ग) $\eta uL/h = \num{3.8e12} +
\num{4.6e12}$ से $u = \num{8.4e12}\times\num{1.5e5}/(\num{6e6}\times10^{20})
= \qty{2.1e-9}{m.s^{-1}} = \qty{6.6}{cm/yr}$ प्राप्त होता है, जो तेज़
महासागरीय प्लेटों की चाल है।
\end{solution}

\begin{solution}{pb:b3:mantle-convection:1}
\textbf{1.} $p = 1000\times9.8\times300 = \qty{2.94}{MPa}$।

\textbf{2.} $w = \rho_wh/\rho_m = 300\times1000/3300 = \qty{90.9}{m}$। भूपर्पटी
एक खंड के रूप में नीचे खिसकती है; भार के नीचे से जो बाहर बहता है और वापस लौटता है,
वह प्रावार है, इसलिए स्तंभ गहराई पर विस्थापित प्रावार से संतुलित होता है।

\textbf{3.} $\pi\times(\num{1.25e5})^2 = \qty{4.91e10}{m^2}$, और जल का
वज़न $\num{4.91e10}\times300\times1000 = \qty{1.47e16}{kg}$ है।

\textbf{4.} तटरेखा समतल बनी थी, जब तल केंद्र पर सबसे अधिक अवनमित था,
तब से केंद्र किनारे की तुलना में अधिक उठ चुका है; इसलिए रेखा अब
केंद्र की ओर चढ़ती है, \qty{125}{km} पर \qty{77.2}{m} ऊँचाई के अंतर से:
\qty{0.62}{m/km}।

\textbf{5.} $\lambda/2 = \qty{250}{km}$, इसलिए $\lambda = \qty{500}{km}$ और $k =
2\pi/\num{5e5} = \qty{1.257e-5}{m^{-1}}$।

\textbf{6.} $1/k = \lambda/2\pi = \qty{79.6}{km}$।

\textbf{7.} $w = 90.9 - \text{उत्थान}$: 79.4, 69.3, 52.9, 40.4, 26.9 और
\qty{13.7}{m}।

\textbf{8.} $w = w_0e^{-t/\tau}$ से $\ln w = \ln w_0 - t/\tau$ प्राप्त होता है,
जिसकी ढलान $-1/\tau$ की एक सीधी रेखा है।

\textbf{9.} $\tau = (14 - 1)/\ln(79.4/13.7) = 13/1.7571 = \qty{7.40}{kyr}$।

\textbf{10.} $w(6) = 79.4\,e^{-5/7.40} = \qty{40.4}{m}$ है, जो
तालिका के अनुसार \qty{50.5}{m} का उत्थान है।

\textbf{11.} \qty{13.7}{m} शेष हैं; $\dot w = 13.7/7400 = \qty{1.85}{mm/yr}$।

\textbf{12.} पूरे \qty{13.7}{m} तक पहुँचता है, पर अधिकाधिक धीमी गति से।

\textbf{13.} दर $e^{-t/\tau}$ के अनुसार घटती है: $t = 7.40\ln(1.85/0.5) =
7.40\times1.309 = \qty{9.7}{kyr}$, लगभग 9700 वर्ष और।

\textbf{14.} $\tau = 7399\times\num{3.156e7} = \qty{2.335e11}{s}$ और $\eta =
\tau\rho_mg\lambda/(4\pi) = \num{2.335e11}\times3300\times9.8\times\num{5e5}/
12.566 = \qty{3.0e20}{Pa.s}$।

\textbf{15.} फेनोस्कैंडियन परास के निचले छोर से ठीक नीचे। \qty{500}{km} तरंगदैर्घ्य के भार के नीचे का प्रवाह
\qty{80}{km} के भीतर क्षीण होता है:
प्रतिक्षेप सबसे उथले, सबसे दुर्बल प्रावार का नमूना लेता है, जो उसके नीचे और कुछ हद तक
उस भाग के भीतर है जो वास्तविक पृथ्वी में स्थलमंडल है।

\textbf{16.} $\tau = 4\pi\times\num{3.0e20}/(3300\times9.8\times\num{3e6}) =
\qty{3.9e10}{s} = \qty{1.2}{kyr}$। इतने चौड़े भार के नीचे प्रवाह
\qty{477}{km} और उससे अधिक गहराई तक पहुँचता है, उस निचले प्रावार की ओर जो दस गुना अधिक श्यान है,
जो उसे धीमा कर देता है।

\textbf{17.} स्थलमंडल संकीर्ण भार का कुछ भाग आनमन से सँभालता है: संतुलन पर अवनमन
समस्थितिक \qty{90.9}{m} की तुलना में कम गहरा और अधिक चौड़ा है, और
प्रतिक्षेप का कुछ भाग प्रत्यास्थ और तात्कालिक है। उत्थान से मेल खाने वाला एकसमान श्यान अर्ध-समष्टि
फिर ऐसी श्यानता देता है जिसमें दोनों प्रभाव मिले होते हैं।

\textbf{18.} $\tau_M = \num{3.0e20}/\num{7.7e10} = \qty{3.9e9}{s}$, लगभग
\qty{124}{yr}, $\tau$ से साठ गुना छोटा: हज़ारों वर्षों में प्रावार
श्यान तरल की भाँति प्रतिक्रिया देता है।

\textbf{19.} $1 - e^{-50/7.40} = 0.9988$: \qty{99.9}{\percent} संतुलन अवनमन का
प्राप्त हो चुका था, और धारणा सही है।

\textbf{20.} $\eta \propto \lambda$: $\lambda$ को आधा करने से $\eta$ भी आधा होता है, और
$\log_{10}\eta$ का मान $\log_{10}2 = 0.30$ से घटता है।

\textbf{21.} $\mathrm{Ra} = 3700\times9.9\times\num{3.1e-5}\times1450\times
(\num{6.7e5})^3/(10^{-6}\times\num{3.0e20}) = \num{1.65e6}$।

\textbf{22.} $\num{1.65e6}/657.5 \approx 2500$ गुना क्रांतिक मान।

\textbf{23.} $(\mathrm{Ra}/2\mathrm{Ra}_c)^{1/3} = (1254)^{1/3} = 10.78$: $\delta
= 670/10.78 = \qty{62}{km}$ और $\mathrm{Nu} = 5.4$।

\textbf{24.} $u = (10^{-6}/\num{6.7e5})\times10.78^2 = \qty{1.73e-10}{m.s^{-1}}
= \qty{0.55}{cm/yr}$, प्लेटों से दो से चौदह गुना धीमा।

\textbf{25.} $q = 5.4\times3.1\times1450/\num{6.7e5} = \qty{36}{mW.m^{-2}}$।
परत प्रावार के भीतर उत्पन्न ऊष्मा और
क्रोड तथा निचले प्रावार की ऊष्मा को छोड़ देती है; और स्केलन कोटि-एक के अपने स्थिरांकों को हटा देता है,
जबकि प्लेटें, कटकों से अर्ध-समष्टि की भाँति ठंडी होती हुई, नवीन तल पर
माध्य सीमा-परत की तुलना में कहीं अधिक ऊष्मा खोती हैं।

\textbf{26.} $\eta = \qty{3.0e20}{Pa.s}$: $\log_{10}\eta =
\textbf{\num[text-series-to-math=true]{20.48}}$।
\end{solution}
